% source: manti.01.ma, 1ère SM, série suites (Drive 0B6pfSJcEVng0bjRLdzhVRjRZY3M).
% Ex. 5 = suites_s1 Ex. 4 ; Ex. 7 = MUFx Ex. 1 ; Ex. 8 = RzBR Ex. 3 (non repris)
\ex{Suites}{bjRL}{Ex. 1}{7}
Soit $(U_n)_{n\in\mathbb{N}^*}$ la suite définie par $U_n=\displaystyle\sum_{k=1}^{n}\frac{1}{n+k}$.
\begin{enumerate}
\item Calculer $U_1$ et $U_3$.
\item Montrer que $(\forall n\in\mathbb{N}^*)\ U_n\geq\dfrac12$.
\item Étudier la monotonie de la suite $(U_n)_{n\in\mathbb{N}^*}$.
\end{enumerate}
\fin
\ex{Suites}{bjRL}{Ex. 2}{7}
Soit $(U_n)_{n\in\mathbb{N}^*}$ la suite définie par $U_n=\displaystyle\sum_{k=1}^{n}\frac{1}{n+k^2}$.
\begin{enumerate}
\item Calculer $U_1$ et $U_3$.
\item Montrer que $(\forall n\in\mathbb{N}^*)\ 0<U_n<1$.
\item Étudier la monotonie de la suite $(U_n)_{n\in\mathbb{N}^*}$.
\end{enumerate}
\fin
% Ex. 3 = série française suites_s1 (non repris)
\ex{Suites}{bjRL}{Ex. 4}{6}
On considère la suite $(U_n)_{n\in\mathbb{N}}$ définie par $U_0=2$ et $U_{n+1}=\sqrt{6+U_n}$.
\begin{enumerate}
\item Calculer $U_1$ et montrer que $(\forall n\in\mathbb{N})\ 2\leq U_n<3$.
\item Étudier la monotonie de la suite $(U_n)_{n\in\mathbb{N}}$.
\end{enumerate}
\fin
\ex{Suites}{bjRL}{Ex. 6}{8}
Soit $(V_n)_{n\geq0}$ la suite définie par $V_0=2$ et $V_{n+1}=\dfrac{7V_n-25}{V_n-3}$.
\begin{enumerate}
\item Montrer que $(\forall n\in\mathbb{N})\ V_n\neq5$.
\item On pose $Y_n=\dfrac{1}{V_n-5}$. Montrer que $(Y_n)_{n\geq0}$ est une suite arithmétique, puis calculer $V_n$ en fonction de $n$.
\item Calculer la somme $Y_0+Y_1+\dots+Y_{10}$.
\end{enumerate}
\fin
\ex{Suites}{bjRL}{Ex. 9}{7}
Soit $(U_n)_{n\in\mathbb{N}}$ la suite définie par $U_0=\dfrac12$ et $U_{n+1}=\dfrac{U_n}{1+U_n}$.
\begin{enumerate}
\item Montrer que $(U_n)_{n\in\mathbb{N}}$ est minorée par $0$.
\item Étudier la monotonie de la suite $(U_n)_{n\in\mathbb{N}}$.
\item On pose $V_n=\dfrac{1}{2U_n}$.
\begin{enumerate}
\item Montrer que $(V_n)_{n\geq0}$ est une suite arithmétique et calculer $V_n$ en fonction de $n$.
\item Calculer $U_n$ en fonction de $n$.
\item Calculer en fonction de $n$ la somme $S=\dfrac{1}{U_0}+\dfrac{1}{U_1}+\dots+\dfrac{1}{U_n}$.
\end{enumerate}
\end{enumerate}
\fin
